\section{行业洞察：学术界 vs 工业界}

\subsection{PhD 还值得读吗？}

翁家翌对 2025 年面临选择的年轻人给出了明确建议：

\begin{importantbox}{翁家翌的职业建议}
\textbf{如果你想进工业界，那么读 PhD 就是浪费生命。}你完全可以用 Master 甚至本科为跳板，通过积累差异化的项目经验（特别是 infra 方面）来与 PhD 候选人同台竞争。"招人最主要的目的是招能用、能干活的人。"
\end{importantbox}

他认为学术界正在被重构。当前 AI Lab 最需要的是 infra 人才——infra 是一个"无底洞"，有 research 直觉的人本来就那么几个（在行业干了三年以上的人屈指可数），剩下的瓶颈全在 infra。

翁家翌建议年轻人的思路应该是：先弄清楚 AI Lab 到底需要什么样的人，然后让自己去匹配。如果他们更需要 infra 的人，就多做 infra 的活，哪怕没有 PhD degree 也没有关系——"更重要的是看你的经验能不能 match、有没有用"。他强调，如果你是一个 new grad（应届毕业生），有一段与目标工作高度匹配的经验，"可以抵好几年的工作经历"。

\begin{knowledgebox}{评价体系的演变}
翁家翌的经历展示了评价体系从学校到工业界的演变：清华用 GPA → 导师用论文/比赛/GitHub star → OpenAI 用 infra 贡献和模型发布次数。每一个阶段的评价标准都不同，关键是\textbf{提前识别下一阶段的评价标准}，而不是在当前阶段过度优化。这本身就是一种"投资未来"的策略。
\end{knowledgebox}

\subsection{学术界 RL 研究与工业界的脱节}

翁家翌指出了学术界和工业界在 RL 领域的巨大鸿沟：

\begin{table}[H]
\centering
\begin{tabular}{ll}
\toprule
\textbf{学术界} & \textbf{工业界} \\
\midrule
对着 Atari / MuJoCo 几个 task overfit & 用 RL 解决真实问题（如 LLM 对齐） \\
比谁在 100K step 时分数高 & 比谁的模型在真实用户场景更好 \\
追求新算法发论文 & 追求 infra 正确性和迭代速度 \\
几块卡的小实验 & 大规模分布式训练 \\
\bottomrule
\end{tabular}
\caption{RL 研究的学术界 vs 工业界对比}
\end{table}

翁家翌在 2022 年 8 月意识到这一点后，逐步停止了天授的开发——因为天授仍然是针对 toy benchmark 的，而他应该投入更有意义的事情。他说当时明确意识到"我应该投入更多的时间到更有意义的事情里面"——一旦看清了学术界 RL 和工业界 RL 的本质区别，继续维护针对 toy benchmark 的框架就失去了吸引力。

翁家翌还引用了他同事的一句话来总结这种差距：

\begin{importantbox}{Idea is cheap --- 验证 idea 的能力才贵}
"Idea 非常便宜。你要做的就是在单位时间内能够验证多少有效的 idea，并且要是正确的 infra、正确的结果、快速的迭代。"动脑子的人可能是像 Alec Radford 那样从 GPT-1 就开始做的人，他的 research 直觉比普通 PhD 动脑子更有用——所以 idea 找他讨论就好了。剩下的瓶颈全在 execution。
\end{importantbox}

\subsection{本章小结}

在 AI 快速发展的今天，传统的学术路径（PhD → Paper → 教职）正在被重新评估。翁家翌的经验表明，工程能力、差异化的项目经验和对需求的精准把握，可能比学历和论文数量更能决定一个人的职业发展。


